% Generated draft: requires mathematical and source review.
\documentclass[11pt]{article}
\usepackage[margin=25mm]{geometry}
\usepackage{fontspec}
\setmainfont{Latin Modern Roman}
\usepackage{amsmath,amssymb,mathtools,unicode-math}
\setmathfont{Latin Modern Math}
\usepackage{enumitem,needspace,xurl,hyperref}
\setlength{\parindent}{0pt}
\setlength{\parskip}{5pt}
\setlength{\emergencystretch}{3em}
\newcommand{\N}{\mathbb{N}}
\newcommand{\Z}{\mathbb{Z}}
\newcommand{\Q}{\mathbb{Q}}
\newcommand{\R}{\mathbb{R}}
\newcommand{\C}{\mathbb{C}}
\newcommand{\F}{\mathbb{F}}
\newcommand{\A}{\mathbb{A}}
\newcommand{\PP}{\mathbb{P}}
\newcommand{\E}{\mathbb{E}}
\newcommand{\eps}{\varepsilon}
\newcommand{\dd}{\,\mathrm{d}}
\begin{document}
\section*{UnsolvedMath ID: 2307008}
\textbf{Source problem code:} catalogue entry
\subsection{Definitions and mathematical statement}
In a standard analytic convention, a quasiconformal map in $\R^3$ is a Sobolev homeomorphism with bounded distortion. To make one common maximal-dilatation convention explicit, put
\[
K_O(f)=\inf\{K\ge1:\|Df\|^3\le K|J_f|\text{ a.e.}\},\qquad
K_I(f)=\inf\{K\ge1:|J_f|\le K\ell(Df)^3\text{ a.e.}\},
\]
\[
K(f)=\max\{K_O(f),K_I(f)\},
\]
where $\|Df\|$ and $\ell(Df)$ are the largest and smallest singular values and $J_f$ is the Jacobian. These inequalities avoid an undefined ratio at a zero derivative. The source does not specify whether maximal dilatation means this two-sided quantity, outer dilatation alone, or another convention; the distinction matters in dimension three.

On the global self-homeomorphism interpretation of ``in $3$-space'', Gehring asks whether
\[
\forall K>1\ \forall f:\R^3\to\R^3\text{ quasiconformal with }K(f)\le K,
\]
\[
\exists g,h:\R^3\to\R^3\text{ quasiconformal},\qquad
f=g\circ h,\qquad K(g)<K,\quad K(h)<K.
\]
The factors may depend on $f$ and $K$. No quantitative uniform improvement below $K$, formula for the factors, or prescribed allocation of distortion is requested. Requiring the strict improvement relative to the actual value $K(f)>1$ is another natural precise phrasing of the same intended factorization issue, but the source does not distinguish exact from upper-bound use of the label $K$.

The condition $K>1$ is necessary for this strict wording: every maximal dilatation is at least one, so the source's unqualified $K=1$ case cannot ask for both factors to have dilatation below one. This endpoint is flagged rather than silently included in an impossible assertion. The extract also leaves open whether the maps are global self-maps or homeomorphisms between arbitrary domains, what intermediate domain is allowed, and whether orientation preservation is required. For a domain version the appropriate formulation would be $h:\Omega\to U$, $g:U\to\Omega'$ with $f=g\circ h$, but neither the domain class nor $U$ is supplied.

The source says the corresponding planar result is true. Its supplied triage labels the three-dimensional question open as of the 2018 edition and reports no definitive later resolution. These are attributed source statements, not independently verified results.
\subsection{Short English statement}
Can every nonconformal quasiconformal map in three dimensions be factored into two maps each with strictly smaller maximal dilatation? The strict comparison requires $K>1$. The source leaves the global-versus-domain setting and dilatation convention unspecified and reports only that the planar analogue is true.
\subsection{Sources}
\begin{itemize}
\item Local source registry: \url{data/unsolvedmath/problems.json}, permanent UnsolvedMath ID $2307008$, source code catalogue entry, including the supplied OPEN-TRIAGE assessment. Dataset attribution: Ulam AI and UnsolvedMath Contributors; original-source attribution and rights are retained.
\item Hayman, \emph{Research Problems in Function Theory}, source-list date 2018, Problem 7.8, attributed to F. W. Gehring: \url{https://arxiv.org/abs/1809.07200}. The metadata reports source-TeX extraction and July 29, 2026 access, with current-status and rights review flags.
\item The statement says the planar analogue holds but supplies no separate theorem citation. No independent source inspection or literature verification was performed for this draft.
\end{itemize}
\end{document}
